% ============================================================================
% 03_optimization | 03_quantization
% Section-body file: include immediately after the relevant \section command
% in the master tstextbook.tex file.
%
% This file intentionally contains no document class, package imports,
% \chapter, \section, or bibliography command. Figure source fragments are
% stored in the root-level figures directory and inserted with \input.
% References are maintained in the root-level bibliography.bib file.
% ============================================================================

The preceding sections established how an objective function defines the
learning problem and how an optimization algorithm uses gradients to update
the parameters of a neural network. These operations explain how a language
model learns its parameters, but they do not by themselves explain whether
the resulting model can be stored and executed within a practical
computational budget. As the number of parameters and the length of
processed sequences increase, the numerical precision used to represent
weights and intermediate tensors becomes an important part of the model's
computational design.

Quantization addresses this problem by representing numerical quantities
with fewer bits. It can reduce parameter storage, the memory required for
intermediate representations, and, when supported by the hardware and
kernels, the cost of arithmetic operations. The reduction in precision,
however, also reduces the number of representable values and therefore
introduces approximation error. The central problem is consequently not
merely to compress a model, but to choose numerical representations that
reduce computational requirements while preserving the information needed
for accurate prediction and stable optimization.

This distinction connects quantization directly to the optimization
process. A pretrained model can be quantized after its training has
finished, or quantization can be simulated during training so that the
model parameters adapt to the numerical constraints that will be present
at inference time. The first strategy is generally described as
post-training quantization (PTQ), whereas the second is quantization-aware
training (QAT). The surveys by \textcite{chen2026llmquantization} and
\textcite{song2026targetcentricqat} provide complementary foundations for
this discussion: the former organizes large language model quantization
across pretraining, fine-tuning, and inference, while the latter analyzes
QAT according to the particular tensors being quantized.

\subsection{Numerical representation and the need for quantization}
\label{subsec:quant-numerical-representation}

A neural network performs its computations using numerical representations
of model parameters and intermediate quantities. Consider a linear
transformation
\begin{equation}
    \mathbf{y}=\mathbf{W}\mathbf{x}+\mathbf{b},
    \label{eq:quant-linear-layer}
\end{equation}
where
\(\mathbf{W}\in\mathbb{R}^{m\times n}\) is a weight matrix,
\(\mathbf{x}\in\mathbb{R}^{n}\) is an input vector, and
\(\mathbf{b}\in\mathbb{R}^{m}\) is a bias vector. The output belongs to
\(\mathbb{R}^{m}\). In a Transformer, the same principle appears in the
linear projections used by attention and feed-forward sublayers, as well
as in the transformations applied to intermediate representations.

In conventional high-precision computation, tensor elements may be
represented using formats such as FP32 or FP16. Lower-precision formats
include FP8 and FP4, while integer representations include INT8 and INT4.
The bit-width \(b\) determines the maximum number of distinct binary
codewords that can represent a scalar:
\begin{equation}
    N_{\mathrm{levels}}=2^b.
    \label{eq:quant-number-levels}
\end{equation}
For an unsigned integer representation, the range is
\(0,\ldots,2^b-1\). A common signed two's-complement range is
\begin{equation}
    q_{\min}=-2^{b-1},
    \qquad
    q_{\max}=2^{b-1}-1.
    \label{eq:quant-signed-range}
\end{equation}
A smaller bit-width therefore provides fewer representable values. The
result is a more compact numerical representation, but also a greater
risk that distinct real values will be mapped to the same discrete level.

The storage implications can be derived directly. Let a model contain
\(N\) parameters, each represented using \(b\) bits. Ignoring scale
factors, zero-points, alignment, and other metadata, the parameter
storage is
\begin{equation}
    M_b=\frac{Nb}{8}\ \text{bytes}.
    \label{eq:quant-model-memory}
\end{equation}
FP16 storage requires \(2N\) bytes, whereas INT4 storage requires
\(N/2\) bytes. The idealized compression ratio relative to FP16 is
therefore
\begin{equation}
    C_{\mathrm{FP16}\rightarrow b}
    =
    \frac{16}{b}.
    \label{eq:quant-compression-ratio}
\end{equation}
For a model with \(10^9\) parameters, FP16 weights require approximately
\(2\) GB when decimal gigabytes are used, while INT4 weights require
approximately \(0.5\) GB before quantization metadata and implementation
overheads are included. This is an idealized weight-storage calculation,
not a prediction of the total memory required by a training or inference
system. The actual memory footprint also depends on activations,
gradients, optimizer states, cached representations, and the storage of
quantization parameters \parencite{chen2026llmquantization}.

The distinction between parameter storage and total execution memory is
important for large language models. During training, weights are only
one part of the memory requirement: activations must be retained or
recomputed for backpropagation, gradients must be represented, and
optimizers may maintain additional state variables. During autoregressive
inference, the key-value (KV) cache grows with the number of retained
tokens. Quantizing only the weights therefore does not necessarily
resolve every memory bottleneck. This motivates treating quantization
according to both the numerical representation and the computational
role of the quantity being represented.

\subsection{Uniform affine quantization}
\label{subsec:quant-uniform-affine}

A fundamental formulation used in neural-network quantization is uniform
affine quantization. It maps a real-valued tensor to a discrete integer
tensor using equally spaced quantization levels. Let
\begin{equation}
    \mathbf{x}\in
    \mathbb{R}^{d_1\times d_2\times\cdots\times d_k}
    \label{eq:quant-real-tensor}
\end{equation}
be a real-valued tensor. Its integer representation is denoted by
\(\mathbf{x}_{\mathrm{int}}\). Uniform affine quantization is determined
by a positive step size \(s\), an integer zero-point \(z\), and the
representable integer interval \([q_{\min},q_{\max}]\). The step size
sets the distance between adjacent reconstruction levels. The
zero-point determines which integer code represents the real value zero.

For a scalar \(x\), the quantization operation is
\begin{equation}
    x_{\mathrm{int}}
    =
    \operatorname{clamp}
    \left(
        \operatorname{round}\left(\frac{x}{s}\right)+z,
        q_{\min},q_{\max}
    \right),
    \label{eq:quant-affine-quantize}
\end{equation}
where \(\operatorname{round}(\cdot)\) denotes rounding to the nearest
integer and \(\operatorname{clamp}(\cdot)\) restricts its argument to the
valid integer interval. In explicit piecewise form,
\begin{equation}
    \operatorname{clamp}(u,q_{\min},q_{\max})
    =
    \begin{cases}
        q_{\min}, & u<q_{\min},\\
        u, & q_{\min}\leq u\leq q_{\max},\\
        q_{\max}, & u>q_{\max}.
    \end{cases}
    \label{eq:quant-clamp}
\end{equation}
The integer code is converted back to an approximation in the real
domain through dequantization:
\begin{equation}
    \widehat{x}=s(x_{\mathrm{int}}-z).
    \label{eq:quant-affine-dequantize}
\end{equation}
Combining the two operations gives the quantize--dequantize mapping
\begin{equation}
    \boxed{
    Q_{s,z}(x)
    =
    s\left[
        \operatorname{clamp}
        \left(
            \operatorname{round}\left(\frac{x}{s}\right)+z,
            q_{\min},q_{\max}
        \right)-z
    \right].
    }
    \label{eq:quant-affine-q}
\end{equation}
The output \(\widehat{x}=Q_{s,z}(x)\) is the reconstructed value used
to approximate the original real-valued quantity. The corresponding
quantization error is
\begin{equation}
    \epsilon=x-\widehat{x}.
    \label{eq:quant-error-definition}
\end{equation}
The mapping and its constituent operations are summarized in
\cref{fig:quant-affine-map}. This quantize--dequantize form is especially
important in QAT, where the forward computation simulates the effect of
low-precision representation while the training implementation may
continue to use floating-point arithmetic
\parencite{song2026targetcentricqat,chen2026llmquantization}.

\begin{figure}[htbp]
    \centering
    % Figure source fragment: uniform affine quantization and dequantization.
% This file is inserted inside a figure environment using \input.
\begin{center}
\setlength{\fboxsep}{6pt}
\begin{tabular}{c}
    \fbox{\parbox[c][0.82cm][c]{0.55\textwidth}{\centering
        \textbf{Real-valued input}\\
        \(x\in\mathbb{R}\)}}\\[0.25em]
    \(\downarrow\;\text{divide by }s\)\\[0.2em]
    \fbox{\parbox[c][0.82cm][c]{0.55\textwidth}{\centering
        \textbf{Scale}\\
        \(x/s\)}}\\[0.25em]
    \(\downarrow\;\text{round to nearest integer}\)\\[0.2em]
    \fbox{\parbox[c][0.82cm][c]{0.55\textwidth}{\centering
        \textbf{Shift and clip}\\
        \(\operatorname{clamp}(\operatorname{round}(x/s)+z,
        q_{\min},q_{\max})\)}}\\[0.25em]
    \(\downarrow\)\\[0.2em]
    \fbox{\parbox[c][0.82cm][c]{0.55\textwidth}{\centering
        \textbf{Integer code}\\
        \(x_{\mathrm{int}}\in\mathbb{Z}\)}}\\[0.25em]
    \(\downarrow\;\text{dequantize: }s(x_{\mathrm{int}}-z)\)\\[0.2em]
    \fbox{\parbox[c][0.82cm][c]{0.55\textwidth}{\centering
        \textbf{Reconstructed value}\\
        \(\widehat{x}\approx x\)}}\\
\end{tabular}
\end{center}

    \caption{Uniform affine quantization and dequantization. A real-valued
    input is scaled, rounded, shifted by the zero-point, and clipped to
    the integer range. Dequantization maps the integer code back to an
    approximation in the real domain.}
    \label{fig:quant-affine-map}
\end{figure}
% Alt text for accessibility register: A vertical flow diagram showing a
% real-valued input x, division by the step size, rounding, addition of the
% zero-point and clipping to the integer range, followed by dequantization
% using the step size and zero-point to obtain the reconstructed value.

\paragraph{Symmetric quantization.}
In symmetric quantization, the zero-point is fixed to zero:
\begin{equation}
    z=0.
    \label{eq:quant-symmetric-zero-point}
\end{equation}
The quantization and dequantization operations consequently simplify to
\begin{align}
    x_{\mathrm{int}}
    &=
    \operatorname{clamp}
    \left(
        \operatorname{round}\left(\frac{x}{s}\right),
        q_{\min},q_{\max}
    \right),
    \label{eq:quant-symmetric-integer}
    \\
    \widehat{x}
    &=
    s x_{\mathrm{int}}.
    \label{eq:quant-symmetric-reconstruction}
\end{align}
For a symmetric real-valued clipping interval, define
\begin{equation}
    \alpha=\max\{|x_{\min}|,|x_{\max}|\}.
    \label{eq:quant-symmetric-alpha}
\end{equation}
A common symmetric step size is
\begin{equation}
    s=\frac{\alpha}{q_{\max}},
    \label{eq:quant-symmetric-scale}
\end{equation}
where \(q_{\max}>0\). In a symmetric signed quantizer, the negative
integer endpoint may be restricted to \(-q_{\max}\), even when the
underlying signed integer data type has one additional negative code.
This convention keeps the quantization interval symmetric around zero.
Symmetric quantization is commonly used for weight tensors because the
zero-point-free formulation simplifies integer arithmetic
\parencite{song2026targetcentricqat,chen2026llmquantization}.

\paragraph{Asymmetric quantization.}
Asymmetric quantization permits a nonzero zero-point and can therefore
represent a real-valued interval that is not centered around zero. Given
a selected clipping interval \([x_{\min},x_{\max}]\) and integer interval
\([q_{\min},q_{\max}]\), the step size is
\begin{equation}
    s=
    \frac{x_{\max}-x_{\min}}
         {q_{\max}-q_{\min}}.
    \label{eq:quant-asymmetric-scale}
\end{equation}
The zero-point is chosen so that real zero is represented as accurately
as possible:
\begin{equation}
    z=
    \operatorname{round}
    \left(q_{\min}-\frac{x_{\min}}{s}\right),
    \label{eq:quant-asymmetric-zero-point}
\end{equation}
and is clipped to the valid integer range when necessary. Asymmetric
quantization can be useful for activation distributions that are shifted
or predominantly nonnegative. The choice between symmetric and
asymmetric quantization is therefore related to the statistical
distribution of the tensor and to the arithmetic supported by the target
implementation \parencite{song2026targetcentricqat,chen2026llmquantization}.

\paragraph{Quantization granularity.}
The number of tensor elements that share the same quantization
parameters determines the quantization granularity. Per-tensor
quantization uses one scale and zero-point for an entire tensor.
Per-channel quantization assigns separate parameters to channels, which
allows different channel distributions to be represented using different
ranges. Per-group quantization partitions a tensor into smaller groups
and assigns quantization parameters to each group. Per-token
quantization applies parameters separately to token representations and
is relevant when activation or KV-cache distributions vary across token
positions. Per-layer terminology is also used when quantization
parameters are shared at the layer level; the exact distinction between
per-layer and per-tensor terminology depends on which tensors are
included in the implementation.

Finer granularity can reduce distortion because the scale and zero-point
adapt to more local distributions. It also requires additional metadata
and can increase indexing, memory, and computational overhead. Coarser
granularity reduces this overhead but may force numerically heterogeneous
values to share a quantization range. Group size, tensor layout, and the
availability of efficient kernels must therefore be considered together
\parencite{song2026targetcentricqat,chen2026llmquantization}.

\subsection{Quantization error and its propagation through neural networks}
\label{subsec:quant-error-propagation}

Quantization error has two principal sources. Rounding error occurs
because a real-valued quantity within the representable interval must
be mapped to one of the discrete quantization levels. Clipping error
occurs when the quantity lies outside the representable interval and
must be saturated at a boundary. With a consistent sign convention,
the total error can be decomposed into rounding and clipping
contributions:
\begin{equation}
    \epsilon=\epsilon_{\mathrm{round}}
              +\epsilon_{\mathrm{clip}}.
    \label{eq:quant-error-decomposition}
\end{equation}
For a value that is not clipped, nearest-level rounding gives the
absolute-error bound
\begin{equation}
    |\epsilon_{\mathrm{round}}|
    \leq \frac{s}{2}.
    \label{eq:quant-rounding-bound}
\end{equation}
This bound does not apply to clipped values, whose error can be much
larger than half a quantization step.

A numerical example makes the distinction explicit. Consider a
symmetric INT4 quantizer with the symmetric integer interval
\([-7,7]\) and real clipping interval \([-1,1]\). The step size is
\begin{equation}
    s=\frac{1}{7}\approx 0.142857.
    \label{eq:quant-example-step}
\end{equation}
For \(x=0.42\), the integer code is
\begin{equation}
    x_{\mathrm{int}}
    =
    \operatorname{round}(0.42/0.142857)
    =
    \operatorname{round}(2.94)=3.
    \label{eq:quant-example-round}
\end{equation}
The reconstructed value is approximately
\begin{equation}
    \widehat{x}=3(0.142857)=0.428571,
    \label{eq:quant-example-reconstruction}
\end{equation}
so the signed error is
\begin{equation}
    \epsilon=0.42-0.428571\approx -0.008571.
    \label{eq:quant-example-round-error}
\end{equation}
This is a rounding error. For \(x=1.2\), the unbounded rounded code
would be \(8\), but the largest permitted code is \(7\). The value is
therefore clipped, giving
\begin{equation}
    x_{\mathrm{int}}=7,\qquad
    \widehat{x}=1,\qquad
    \epsilon=1.2-1=0.2.
    \label{eq:quant-example-clipping-error}
\end{equation}
The example illustrates the trade-off in selecting a clipping range:
expanding the range reduces clipping but increases the spacing between
quantization levels, whereas narrowing the range improves resolution
for values near zero but can clip more extreme values
\parencite{song2026targetcentricqat,chen2026llmquantization}.

Quantization errors also propagate through the network. Consider the
linear transformation in \cref{eq:quant-linear-layer}. Let the
quantized weight matrix and input vector be
\begin{align}
    \widehat{\mathbf{W}}
    &=\mathbf{W}+\Delta\mathbf{W},
    \label{eq:quant-weight-perturbation}
    \\
    \widehat{\mathbf{x}}
    &=\mathbf{x}+\Delta\mathbf{x},
    \label{eq:quant-activation-perturbation}
\end{align}
where \(\Delta\mathbf{W}\) and \(\Delta\mathbf{x}\) are their
quantization perturbations. The quantized output is
\begin{equation}
    \widehat{\mathbf{y}}
    =
    (\mathbf{W}+\Delta\mathbf{W})
    (\mathbf{x}+\Delta\mathbf{x})+\mathbf{b}.
    \label{eq:quant-perturbed-output}
\end{equation}
Expanding the product and subtracting the original output gives
\begin{equation}
    \boxed{
    \Delta\mathbf{y}
    =
    \Delta\mathbf{W}\mathbf{x}
    +\mathbf{W}\Delta\mathbf{x}
    +\Delta\mathbf{W}\Delta\mathbf{x}.
    }
    \label{eq:quant-output-error}
\end{equation}
If only weights are quantized, then
\(\Delta\mathbf{x}=\mathbf{0}\), and
\begin{equation}
    \Delta\mathbf{y}
    =
    \Delta\mathbf{W}\mathbf{x}.
    \label{eq:quant-weight-only-error}
\end{equation}
If only activations are quantized, then
\(\Delta\mathbf{W}=\mathbf{0}\), and
\begin{equation}
    \Delta\mathbf{y}
    =
    \mathbf{W}\Delta\mathbf{x}.
    \label{eq:quant-activation-only-error}
\end{equation}
When both are quantized, the interaction term
\(\Delta\mathbf{W}\Delta\mathbf{x}\) also contributes. This is one
reason joint weight--activation quantization is more complicated than
weight-only quantization.

In a deep network, perturbations can pass through successive linear
transformations, nonlinear activation functions, attention operations,
normalization, and residual connections. The error at a later layer
therefore depends on both the local approximation and the transformations
applied to the perturbed representation. Weight errors are relatively
static after offline quantization, whereas activation errors depend on
the input and runtime tensor distributions. The LLM quantization survey
describes this distinction as static weight error and dynamic activation
error, with activation errors potentially accumulating across layers
\parencite{chen2026llmquantization}. The QAT survey further emphasizes
that error characteristics depend on the quantization target and its
lifetime, reuse frequency, and deployment constraints
\parencite{song2026targetcentricqat}.

\subsection{Quantization as part of the optimization process}
\label{subsec:quantization-optimization}

\paragraph{Post-training quantization.}
Post-training quantization applies low-precision representations to a
model after high-precision training has been completed. The pretrained
parameters are analyzed, quantization parameters are selected, and the
model tensors are mapped to their lower-precision representations.
Calibration data can be used to estimate activation ranges, while
reconstruction objectives or sensitivity information can be used to
choose quantization parameters. PTQ avoids repeating the original
pretraining process and is therefore attractive when training the model
again would be prohibitively expensive.

Its principal limitation is that the original model was optimized
without explicitly accounting for the approximation errors introduced
by quantization. A low-bit representation can consequently alter
sensitive weights or activations and degrade the model's performance.
This issue is especially important when activation distributions contain
outliers: a small number of large values can expand the quantization
range and leave fewer effective levels for the majority of values.
PTQ methods therefore differ in how they estimate ranges, reconstruct
weights, use calibration data, and account for sensitivity
\parencite{chen2026llmquantization}.

\paragraph{Quantization-aware training.}
Quantization-aware training incorporates quantization effects into the
training or fine-tuning computation. Instead of optimizing a model only
under high-precision forward computation, QAT inserts fake-quantization
operations so that the forward pass reflects the approximation that
will be present in the deployed model. If a layer is
\begin{equation}
    \mathbf{y}=f(\mathbf{W}\mathbf{x}+\mathbf{b}),
    \label{eq:quant-full-precision-layer}
\end{equation}
a QAT forward pass can use
\begin{equation}
    \widehat{\mathbf{y}}
    =
    f\left(
        Q(\mathbf{W})Q(\mathbf{x})+\mathbf{b}
    \right),
    \label{eq:quant-qat-forward}
\end{equation}
where \(Q(\cdot)\) denotes the quantize--dequantize operator and only
the selected tensors are quantized. The forward computation is commonly
performed in floating-point arithmetic while the low-precision effects
are emulated by the fake quantizer.

The underlying trainable parameters can remain in full precision as
master weights. Their fake-quantized versions are used in the forward
pass, while the optimizer updates the master parameters. A basic
gradient-descent update is
\begin{equation}
    \mathbf{W}_{t+1}
    =
    \mathbf{W}_t
    -
    \eta
    \nabla_{\mathbf{W}}\mathcal{L},
    \label{eq:quant-master-weight-update}
\end{equation}
where \(\eta\) is the learning rate and \(\mathcal{L}\) is the training
loss. The updated master weights are quantized again during the next
forward pass. Maintaining continuous master weights preserves small
parameter updates that could be lost if optimization were performed
directly on a coarse discrete grid. The distinction between the
full-precision parameter state and its fake-quantized forward proxy is
central to common QAT formulations
\parencite{song2026targetcentricqat}.

\begin{figure}[htbp]
    \centering
    % Figure source fragment: simplified quantization-aware training loop.
% This file is inserted inside a figure environment using \input.
\begin{center}
\setlength{\fboxsep}{6pt}
\begin{tabular}{c}
    \fbox{\parbox[c][0.82cm][c]{0.60\textwidth}{\centering
        \textbf{Full-precision master weights}\\
        \(\mathbf{W}_t\)}}\\[0.25em]
    \(\downarrow\;\text{fake quantization }Q(\cdot)\)\\[0.2em]
    \fbox{\parbox[c][0.82cm][c]{0.60\textwidth}{\centering
        \textbf{Quantized forward proxy}\\
        \(\widehat{\mathbf{W}}_t=Q(\mathbf{W}_t)\)}}\\[0.25em]
    \(\downarrow\;\text{forward computation with quantized proxies}\)\\[0.2em]
    \fbox{\parbox[c][0.82cm][c]{0.60\textwidth}{\centering
        \textbf{Loss evaluation}\\
        \(\mathcal{L}(\widehat{\mathbf{W}}_t,\widehat{\mathbf{x}})\)}}\\[0.25em]
    \(\downarrow\;\text{backpropagation with a surrogate gradient}\)\\[0.2em]
    \fbox{\parbox[c][0.82cm][c]{0.60\textwidth}{\centering
        \textbf{Optimizer update}\\
        \(\mathbf{W}_{t+1}=\mathbf{W}_t-\eta\widetilde{\nabla}_{\mathbf{W}}\mathcal{L}\)}}\\[0.25em]
    \(\downarrow\;\text{repeat the next training step}\)\\[0.2em]
    \(\boxed{\text{Updated master weights}}\)
\end{tabular}
\end{center}

    \caption{Simplified quantization-aware training loop. Master weights
    are fake-quantized for the forward computation; the loss is
    differentiated using a surrogate gradient, and the optimizer updates
    the underlying master weights.}
    \label{fig:quant-qat-flow}
\end{figure}
% Alt text for accessibility register: A loop showing full-precision
% master weights passing through fake quantization into a quantized
% forward computation, followed by loss evaluation, surrogate-gradient
% backpropagation, and an optimizer update to the master weights.

\paragraph{The straight-through estimator.}
The central mathematical difficulty in QAT is that rounding is
non-differentiable. The function
\begin{equation}
    q(x)=\operatorname{round}(x)
    \label{eq:quant-round-function}
\end{equation}
is piecewise constant. Its derivative is zero almost everywhere and
undefined at rounding thresholds. Direct differentiation through the
rounding operation would therefore fail to provide useful gradients
for updating the underlying continuous parameters.

The straight-through estimator (STE) replaces the true derivative
with a surrogate derivative. For a quantized scalar with clipping
interval \([x_{\min},x_{\max}]\), a common clipped STE approximation is
\begin{equation}
    \frac{\partial\widehat{x}}{\partial x}
    \approx
    \begin{cases}
        1, & x_{\min}\leq x\leq x_{\max},\\
        0, & \text{otherwise}.
    \end{cases}
    \label{eq:quant-clipped-ste}
\end{equation}
Using the chain rule with this approximation gives
\begin{equation}
    \frac{\partial\mathcal{L}}{\partial x}
    \approx
    \frac{\partial\mathcal{L}}{\partial\widehat{x}}
    \mathbb{I}(x_{\min}\leq x\leq x_{\max}),
    \label{eq:quant-ste-gradient}
\end{equation}
where \(\mathbb{I}(\cdot)\) is the indicator function. Inside the
clipping interval, the quantizer is treated as if it had the derivative
of the identity function. Outside the interval, the gradient is masked.
This is a surrogate-gradient rule, not the true derivative of the
discrete quantization operator. It makes gradient-based optimization
practical but introduces a mismatch between the forward quantizer and
the backward approximation \parencite{song2026targetcentricqat}.

The clipped STE also has a limitation: saturated values receive zero
gradient through the quantizer and may consequently have difficulty
returning to the valid range. This motivates methods that learn or
adapt clipping thresholds, scales, and offsets. PACT learns activation
clipping thresholds, LSQ treats the step size as a learnable parameter,
and LSQ+ introduces learnable offsets for asymmetric activation
distributions. Gradient scaling and magnitude-aware differentiation
address the scale of updates to quantization parameters, while
oscillation-aware methods attempt to reduce repeated threshold
crossings. Differentiable soft quantization replaces hard rounding
during some or all of training with a smooth relaxation, but a
relaxation can introduce a discrepancy between the training operator
and the final hard quantizer. These alternatives are reviewed in
\textcite{song2026targetcentricqat}.

\paragraph{Learning the quantization step size.}
A fixed quantization grid may not remain suitable as the weight and
activation distributions change during training. Learnable step-size
quantization treats the scale \(s\) as an optimization variable. For
symmetric quantization, the fake-quantized scalar is
\begin{equation}
    \widehat{x}
    =
    s\,
    \operatorname{clamp}
    \left(
        \operatorname{round}\left(\frac{x}{s}\right),
        q_{\min},q_{\max}
    \right).
    \label{eq:quant-learnable-step-output}
\end{equation}
Let \(v=x/s\). Under the STE approximation, the derivative with respect
to \(s\) is represented piecewise as
\begin{equation}
    \frac{\partial\widehat{x}}{\partial s}
    \approx
    \begin{cases}
        q_{\min}, & v\leq q_{\min},\\
        \operatorname{round}(v)-v,
            & q_{\min}<v<q_{\max},\\
        q_{\max}, & v\geq q_{\max}.
    \end{cases}
    \label{eq:quant-lsq-step-gradient}
\end{equation}
The loss gradient is then
\begin{equation}
    \frac{\partial\mathcal{L}}{\partial s}
    =
    \frac{\partial\mathcal{L}}{\partial\widehat{x}}
    \frac{\partial\widehat{x}}{\partial s}.
    \label{eq:quant-lsq-chain-rule}
\end{equation}
The learned step size changes the trade-off between clipping-induced
distortion and rounding-induced error. LSQ additionally scales the
gradient of the step-size parameter to make its update magnitude
compatible with updates to the model weights
\parencite{song2026targetcentricqat}.

\paragraph{Quantization oscillation and optimization stability.}
The forward quantizer is piecewise constant, so small changes in a
full-precision parameter can cause its quantized value to jump between
adjacent levels. For a uniform quantizer with step size \(s\), a
rounding threshold between levels \(k\) and \(k+1\) is
\begin{equation}
    w_{\mathrm{threshold}}
    =
    \left(k+\frac{1}{2}\right)s.
    \label{eq:quant-threshold}
\end{equation}
A parameter moving back and forth across this threshold can repeatedly
change its quantized value even when its continuous updates are small.
This behavior is called quantization oscillation and can impair
convergence, particularly at low bit-widths. Oscillation-aware
optimization, adaptive quantization parameters, staged precision
changes, and selected weight-freezing strategies are among the
stabilization directions discussed in
\textcite{song2026targetcentricqat}.

The optimization sensitivity of parameters also explains why equal
quantization errors need not produce equal changes in model loss. Let
\(\mathbf{w}\) be the parameter vector and
\(\Delta\mathbf{w}\) its quantization perturbation. A second-order
Taylor approximation gives
\begin{equation}
    \mathcal{L}(\mathbf{w}+\Delta\mathbf{w})
    \approx
    \mathcal{L}(\mathbf{w})
    +
    \nabla\mathcal{L}(\mathbf{w})^{\top}\Delta\mathbf{w}
    +
    \frac{1}{2}
    \Delta\mathbf{w}^{\top}
    \mathbf{H}
    \Delta\mathbf{w},
    \label{eq:quant-loss-taylor}
\end{equation}
where \(\mathbf{H}=\nabla^2\mathcal{L}(\mathbf{w})\) is the Hessian.
Near a stationary point, the gradient term may be small, and the
quadratic term describes how the curvature of the loss surface weights
different perturbation directions. This provides a mathematical
motivation for sensitivity-aware reconstruction and mixed-precision
allocation. The Hessian expression is a local approximation; it does
not by itself specify a complete quantization algorithm.

\subsection{Quantization targets and the lifecycle of a language model}
\label{subsec:quant-targets-lifecycle}

Quantization can be applied at different stages of the language-model
lifecycle. The target-centric perspective is important because weights,
activations, gradients, and cached representations have different
statistical properties, lifetimes, reuse frequencies, and computational
roles. \textcite{song2026targetcentricqat} organizes QAT around weights,
activations, joint weight--activation quantization, KV cache, and
gradients. \textcite{chen2026llmquantization} additionally organizes
large language model quantization around pretraining, fine-tuning, and
inference.

\begin{figure}[htbp]
    \centering
    % Figure source fragment: principal quantization targets in an LLM.
% This file is inserted inside a figure environment using \input.
\begin{center}
\setlength{\fboxsep}{6pt}
\begin{tabular}{ccc}
    & \fbox{\parbox[c][0.85cm][c]{0.36\textwidth}{\centering
        \textbf{Large language model}\\
        Transformer computation}} & \\[0.35em]
    \(\swarrow\) & \(\downarrow\) & \(\searrow\) \\[0.2em]
    \fbox{\parbox[c][0.95cm][c]{0.27\textwidth}{\centering
        \textbf{Weights}\\
        Learned parameters}}
    &
    \fbox{\parbox[c][0.95cm][c]{0.27\textwidth}{\centering
        \textbf{Activations}\\
        Runtime representations}}
    &
    \fbox{\parbox[c][0.95cm][c]{0.27\textwidth}{\centering
        \textbf{Weight--activation}\\
        Joint quantization}}
    \\[0.45em]
    \(\searrow\) & \(\downarrow\) & \(\swarrow\) \\[0.2em]
    \multicolumn{3}{c}{
    \fbox{\parbox[c][0.90cm][c]{0.48\textwidth}{\centering
        \textbf{Training and inference targets}}}}\\[0.35em]
    \(\swarrow\) & & \(\searrow\) \\[0.2em]
    \fbox{\parbox[c][0.95cm][c]{0.30\textwidth}{\centering
        \textbf{Gradients}\\
        Optimization updates}}
    &&
    \fbox{\parbox[c][0.95cm][c]{0.30\textwidth}{\centering
        \textbf{KV cache}\\
        Reused token states}}
\end{tabular}
\end{center}

    \caption{Principal quantization targets in large language models.
    Weights and activations participate in the forward computation;
    gradients and optimizer-related quantities arise during training;
    and the key-value cache stores representations reused during
    autoregressive decoding. The categories identify different
    quantization targets rather than mutually exclusive model stages.}
    \label{fig:quant-targets}
\end{figure}
% Alt text for accessibility register: A central large-language-model
% computation box connected to five labeled target boxes: weights,
% activations, joint weights and activations, key-value cache, and
% gradients. Training and inference roles are indicated for each target.

\paragraph{Weight quantization.}
Weights are learned parameters and remain relatively stationary during
inference. They are therefore natural targets for offline quantization.
If the quantized weight matrix is
\begin{equation}
    \widehat{\mathbf{W}}
    =
    \mathbf{W}+\Delta\mathbf{W},
    \label{eq:quant-weight-error-model}
\end{equation}
then the error in a linear layer is
\begin{equation}
    \Delta\mathbf{y}
    =
    \Delta\mathbf{W}\mathbf{x}.
    \label{eq:quant-weight-error-output}
\end{equation}
The effect of weight error depends on the activation vector with which
the weights are multiplied. Weight-only quantization can reduce model
storage and memory traffic without requiring the activations to use the
same low-precision format.

GPTQ and AWQ are representative weight-only PTQ methods discussed by
\textcite{chen2026llmquantization}. GPTQ uses second-order information
to guide weight reconstruction, whereas AWQ uses activation information
to identify salient weight channels and guide scaling. Group-wise
quantization assigns separate scales to smaller groups of weights,
allowing local distributions to be represented more closely than with
a single global scale. At very low bit-widths, selective preservation
of sensitive weights or channels can reduce the loss of model
capability \parencite{chen2026llmquantization,song2026targetcentricqat}.

\paragraph{Activation quantization.}
Activations are generated dynamically from the input data. Their
distributions can vary across layers, examples, and token positions.
This makes their quantization ranges more difficult to determine than
those of fixed inference weights. Outliers are particularly important:
a small number of large values can increase the scale of a symmetric
quantizer and reduce the effective resolution available for most
values. Learnable clipping, per-token or per-channel scales, and
mixed-precision handling of outliers are strategies for addressing
this issue. SmoothQuant and RPTQ are representative weight--activation
quantization methods discussed in the LLM quantization survey
\parencite{chen2026llmquantization}.

\paragraph{Joint weight--activation quantization.}
A notation such as W4A4 indicates that both weights and activations
are represented using four bits. Joint quantization can reduce memory
traffic and enable low-precision matrix multiplication, but the two
error sources interact through
\cref{eq:quant-output-error}. Moreover, quantizing upstream weights
can change the distributions of downstream activations. Consequently,
calibration ranges and quantization parameters may need to account for
the quantized computational graph rather than being selected
independently for each tensor. Integer-compatible execution also
requires attention to nonlinear and normalization operations such as
Softmax, GELU, and LayerNorm. A model with quantized linear layers is
not necessarily an integer-only model if other operations still require
floating-point execution \parencite{chen2026llmquantization}.

\paragraph{Key-value cache quantization.}
Autoregressive language models generate tokens sequentially. During
self-attention, key and value representations for previously processed
tokens are retained so that later decoding steps can reuse them.
Let \(L\) denote the number of Transformer layers, \(B\) the batch
size, \(T\) the number of cached tokens, \(H_{\mathrm{KV}}\) the
number of key-value heads, and \(d_h\) the dimension of each head.
The number of scalar values stored in the key-value cache is
approximately
\begin{equation}
    N_{\mathrm{KV}}
    =
    2LBT H_{\mathrm{KV}}d_h,
    \label{eq:quant-kv-elements}
\end{equation}
where the factor of two accounts for keys and values. If each scalar
uses \(b\) bits, the idealized cache storage is
\begin{equation}
    M_{\mathrm{KV}}
    =
    \frac{2LBT H_{\mathrm{KV}}d_h b}{8}
    \ \text{bytes}.
    \label{eq:quant-kv-memory}
\end{equation}
The cache therefore grows linearly with sequence length and batch size.
Quantization can reduce this memory requirement, but its error has a
distinctive consequence: cached representations are reused in
subsequent attention computations, so an approximation can affect
multiple future decoding steps. Per-token, per-channel, adaptive, and
non-uniform approaches are discussed as ways of accommodating
variation in cached key and value distributions
\parencite{song2026targetcentricqat,chen2026llmquantization}.

\paragraph{Gradient quantization.}
Gradients determine the updates applied to model parameters. Let
\(\mathbf{g}_t=\nabla_{\mathbf{w}}\mathcal{L}\) be the full-precision
gradient and let its quantized approximation be
\begin{equation}
    \widehat{\mathbf{g}}_t
    =
    \mathbf{g}_t+\boldsymbol{\epsilon}_{g,t}.
    \label{eq:quant-gradient-error}
\end{equation}
A gradient-descent update using the quantized gradient is
\begin{equation}
    \mathbf{w}_{t+1}
    =
    \mathbf{w}_t
    -
    \eta\left(
        \mathbf{g}_t+\boldsymbol{\epsilon}_{g,t}
    \right).
    \label{eq:quant-gradient-update}
\end{equation}
The update perturbation is therefore
\begin{equation}
    \Delta\mathbf{w}_{t+1}
    =
    -\eta\boldsymbol{\epsilon}_{g,t}.
    \label{eq:quant-gradient-update-error}
\end{equation}
Gradient quantization can change both the magnitude and direction of
the update. Its evaluation should therefore consider gradient norms,
directional agreement, and the stability of the optimization
trajectory, rather than tensor reconstruction error alone
\parencite{song2026targetcentricqat}.

\paragraph{Quantization during pretraining and fine-tuning.}
During pretraining, low-precision computation can be applied to
forward activations, weights, gradients, and communication quantities,
while selected persistent states retain higher precision. FP8 formats
such as E4M3 and E5M2 allocate bits differently between exponent and
mantissa: E4M3 provides more mantissa precision over a more restricted
range, whereas E5M2 provides a larger exponent range. This distinction
is relevant because weights, activations, and gradients can have
different magnitude distributions. Dynamic scaling and block-wise
scaling are used to adapt the numerical range to tensor or block
statistics. The LLM quantization survey discusses FP8, INT8, INT4,
and FP4 training directions, as well as mixed-precision training
systems and low-bit pretraining
\parencite{chen2026llmquantization}.

Fine-tuning quantization includes PTQ and QAT. PTQ applies a low-bit
representation to an existing pretrained model, while QAT simulates
quantization during fine-tuning so that the model can adapt to the
deployment representation. QAT can be combined with distillation or
parameter-efficient adaptation. LLM-QAT, QA-LoRA, LoftQ, and QEFT
are examples of approaches discussed in the supplied surveys. Their
detailed relationship to parameter-efficient fine-tuning is developed
in the subsequent PEFT section; here, their relevance is that
quantization can be integrated with adaptation rather than treated
only as a final compression step
\parencite{song2026targetcentricqat,chen2026llmquantization}.

\subsection{Numerical formats, mixed precision, and hardware efficiency}
\label{subsec:quant-formats-hardware}

Integer and floating-point formats distribute their representational
capacity differently. Integer quantization uses a discrete range with
a fixed step size over the selected interval. Floating-point formats
represent values through sign, exponent, and mantissa fields, so the
spacing between representable values varies with magnitude. FP8 and
FP4 formats are therefore not interchangeable with INT8 and INT4
merely because their nominal bit-widths are equal. The suitable format
depends on the dynamic range and precision requirements of the target
tensor and on the arithmetic supported by the hardware
\parencite{chen2026llmquantization}.

Non-uniform quantization uses unequally spaced representational levels.
It can allocate more levels to regions where values are concentrated
or where approximation is more important. NF4 is an example discussed
in the LLM quantization survey for representing weights whose
distribution is approximately normal. Such a representation may
reduce distortion for the intended distribution, although its
implementation and hardware support must also be considered
\parencite{chen2026llmquantization}.

Mixed-precision quantization assigns different bit-widths to different
model components. Suppose a model is divided into \(K\) parameter
groups, where group \(k\) contains \(N_k\) parameters and uses \(b_k\)
bits per parameter. Ignoring metadata, the total storage is
\begin{equation}
    M
    =
    \frac{1}{8}
    \sum_{k=1}^{K}N_kb_k.
    \label{eq:quant-mixed-precision-memory}
\end{equation}
A simplified precision-allocation problem can be written as
\begin{equation}
    \begin{aligned}
        \min_{\{b_k\}}\quad&
        \mathcal{L}_{\mathrm{quantized}}(\{b_k\})\\
        \text{subject to}\quad&
        \sum_{k=1}^{K}N_kb_k\leq B,
    \end{aligned}
    \label{eq:quant-bit-allocation}
\end{equation}
where \(B\) is a bit budget. This formulation expresses the objective
of minimizing quantization-induced task degradation under a storage
constraint. In practice, the loss as a function of bit allocation is
not generally known in closed form. Sensitivity estimates, local
reconstruction error, activation statistics, and hardware constraints
are therefore used to guide allocation
\parencite{song2026targetcentricqat,chen2026llmquantization}.

The practical benefit of quantization depends on whether a workload is
limited by data movement or arithmetic throughput. A memory-bound
workload is constrained primarily by moving weights or activations
between memory and compute units. A compute-bound workload is
constrained primarily by the rate of arithmetic operations. Weight-only
quantization can reduce the amount of weight data transferred, whereas
joint weight--activation quantization can additionally exploit
low-precision arithmetic units when suitable kernels are available.
A reduction in bit-width does not, by itself, guarantee a proportional
reduction in latency: dequantization, memory access patterns, kernel
availability, operator fallbacks, and hardware architecture all affect
the realized result \parencite{chen2026llmquantization}.

Accumulator precision is another constraint. In an integer matrix
multiplication, an output element is an accumulated sum:
\begin{equation}
    y_{ij}
    =
    \sum_{k=1}^{n}w_{ik}x_{kj}.
    \label{eq:quant-integer-accumulation}
\end{equation}
Even if the operands use low-bit integers, the sum can require a
wider representation. For example, if symmetric INT4 operands lie
between \(-7\) and \(7\), the maximum magnitude of one product is
\(49\). An accumulation of \(4096\) such maximum positive products
would be
\begin{equation}
    4096\times49=200{,}704.
    \label{eq:quant-accumulator-example}
\end{equation}
This value exceeds the maximum positive value of a signed 18-bit
integer, \(2^{17}-1=131{,}071\), but is within the maximum positive
value of a signed 19-bit integer, \(2^{18}-1=262{,}143\). This
worst-case calculation illustrates why operand precision and
accumulator precision must be considered separately. Actual
implementations choose accumulator formats according to their
supported arithmetic and overflow constraints
\parencite{song2026targetcentricqat,chen2026llmquantization}.

Hardware--algorithm co-design is consequently part of quantization.
A low-bit representation must be supported by kernels that efficiently
load data, apply scales, perform arithmetic, accumulate products, and
write outputs. A weight-only kernel may reduce memory traffic but still
incur dequantization overhead. A weight--activation quantized model
may require low-precision matrix multiplication and compatible
implementations of nonlinear and normalization operations. Hand-tuned
kernels, compiler-based generation, and mixed-precision execution
therefore influence whether a theoretical compression benefit becomes
a practical speed or energy benefit
\parencite{chen2026llmquantization}.

\subsection{Evaluation, limitations, and connection to the next topic}
\label{subsec:quant-evaluation-transition}

The evaluation of a quantized model must distinguish numerical
reconstruction, preservation of model behavior, and deployment
efficiency. A common tensor-level measure is normalized mean squared
error. For an original tensor \(\mathbf{X}\) and its quantized
reconstruction \(Q(\mathbf{X})\), it is
\begin{equation}
    \operatorname{nMSE}(\mathbf{X})
    =
    \frac{
        \|\mathbf{X}-Q(\mathbf{X})\|_F^2
    }{
        \|\mathbf{X}\|_F^2
    },
    \label{eq:quant-nmse}
\end{equation}
where \(\|\cdot\|_F\) is the Frobenius norm. For a matrix
\(\mathbf{A}\),
\begin{equation}
    \|\mathbf{A}\|_F
    =
    \sqrt{\sum_{i,j}a_{ij}^2}.
    \label{eq:quant-frobenius}
\end{equation}
A smaller nMSE indicates a closer tensor reconstruction. It should be
complemented by cosine similarity, clipping and saturation ratios,
and layer-wise error profiles because a single aggregate error can
conceal distortion concentrated in a small number of sensitive
modules \parencite{song2026targetcentricqat}.

For Transformer models, local tensor error should be paired with
measurements of block-output drift, hidden-state similarity, and
divergence between the output distributions of the full-precision
and quantized models. One distributional measure is the
Kullback--Leibler divergence:
\begin{equation}
    D_{\mathrm{KL}}(P\|Q)
    =
    \sum_i P(i)\log\frac{P(i)}{Q(i)},
    \label{eq:quant-kl}
\end{equation}
where \(P\) and \(Q\) are probability distributions. In this setting,
they can represent the token distributions produced by the original
and quantized models on the same input. Cross-entropy and perplexity
also measure the effect of quantization on language modeling. For a
sequence of \(T\) tokens, the average negative log-likelihood is
\begin{equation}
    \mathcal{L}_{\mathrm{NLL}}
    =
    -\frac{1}{T}
    \sum_{t=1}^{T}
    \log P(x_t\mid x_{<t}),
    \label{eq:quant-nll}
\end{equation}
and perplexity is
\begin{equation}
    \operatorname{PPL}
    =
    \exp(\mathcal{L}_{\mathrm{NLL}}).
    \label{eq:quant-perplexity}
\end{equation}
Perplexity is useful for language modeling but does not, by itself,
establish that reasoning, mathematical problem solving, instruction
following, or code generation capabilities have been preserved.

The QAT survey reports that LLM evaluation uses language-modeling
datasets and task benchmarks that include WikiText-2, C4, MMLU,
ARC, GSM8K, MATH or BBH, and HumanEval. These evaluate different
aspects of model behavior, and results should be compared only with
care because the evaluated backbones, training budgets, prompts,
decoding settings, and benchmark protocols can differ. The survey
also notes that reasoning and code-generation evaluation is less
consistently reported than conventional language-modeling and
zero-shot evaluation in the methods it reviews
\parencite{song2026targetcentricqat}.

Deployment evaluation must report the actual memory footprint,
latency, throughput, and hardware conditions. Prefill and autoregressive
decoding should be distinguished because their computational and
memory behavior differs. The nominal bit-width is not a substitute
for measured efficiency: two models with the same weight bit-width
may use different activation formats, group sizes, accumulators,
kernels, or floating-point fallbacks. QAT-specific stability
diagnostics can additionally include gradient cosine similarity,
gradient sign agreement, gradient-norm drift, and the fraction of
weights whose quantized values repeatedly change across adjacent
updates \parencite{song2026targetcentricqat,chen2026llmquantization}.

Quantization nevertheless involves unavoidable trade-offs. Lower
bit-widths reduce storage and can enable more efficient computation,
but they reduce representational capacity and can increase rounding,
clipping, and propagation errors. Fine-grained scales can improve
local approximation at the cost of metadata and implementation
overhead. QAT can recover performance by adapting a model to
quantization constraints, but its surrogate gradients are approximate
and its optimization can be unstable. Hardware support determines
whether the numerical savings translate into real latency, throughput,
or energy improvements.

The central result is that quantization is not a single compression
operation. It is a design problem involving numerical representation,
statistical distributions, optimization dynamics, model components,
and hardware execution. PTQ and QAT provide different ways of
introducing low-precision representations; target-specific methods
address the distinct behavior of weights, activations, gradients,
and KV-cache entries; and mixed-precision allocation balances
accuracy against a resource budget.

This discussion also establishes a connection to the next optimization
topic. Quantization changes how numerical information is represented,
whereas knowledge distillation transfers predictive information from
a teacher model to a student model. The two can be combined, as the
QAT literature illustrates, but they solve different parts of the
problem. The next section therefore develops distillation as a
complementary mechanism for transferring model behavior under
computational and representational constraints.
